\subsubsection{FLISR Edge solution}
\label{sec:flisr_edge}

\noindent As described previously, the FLISR Edge solution proposed consists of two applications of the FLISR concept, one of which is a centralised implementation deployed in a remote location, such
as the DSOs own systems, with the second deployed on LBFM and Recloser assets by equipping them with an additional edge device, running the distributed algorithm, which can act independently if necessary.
This work expands on the previous related research by vastly expanding the capabilities of the FLISR algorithm by introducing bi-directional communication to directly control switching assets, in constrast
to the somewhat limited previous iteration which simply generated an email forwarded to DSO personnel \cite{10.12688/openreseurope.14115.2}. In addition, FLISR Edge introduces the distributed element deployable on a low-cost edge device.
\\
Such an edge device would require compatible interfaces to link with existing equipment on the pole, sufficient processing capability to run Dockerised applications and finally cellular connectivity to
send and receive data and control signals. With the rising popularity of low-cost and Edge computing, coupled with the growing rate of processing available to even low-power chips, the options available
to provision a low-cost Edge computing device is much higher than in the previous decade \cite{8016573}. For example, ``Raspberry Pi'' provide a range of hardware configurations
for low-cost computing, including a wide array of additional modules available to meet a variety of use cases. A review of available retailers suggests that
on average acquiring a low-cost computing device with the appropriate processing capacity (e.g., Raspberry Pi 3 Model B+) and required modules would cost approximately one hundred euros per device minus installation costs.
When considering a network such as Waterford, the cost of retrofitting existing LBFM switches with Reclosers to enable loop automation, costs 61,900 euros per unit according to ESB Networks 2023 standard prices report \cite{Networks_2023}.
In comparison then, exploring the application of low-cost Edge computing devices to enable loop automation on existing assets appears more financially feasible.
\\
\\
Assets equipped with switching hardware consist of a number of modules which allow for interfacing with the DSOs SCADA system for interpreting, sending and receiving data and enacting control operations
on the physical circuit breaker. At the core of this architecture is a Remote Terminal Unit (RTU), a microprocessor based device, which streams real-time status data to and from the connected circuit breaker
\cite{6032722}.
Communication with the DSOs SCADA system conforms to the IEC  60870-5-101/104 protocol, which provides network access based on TCP/IP communications and enables basic control tasks between centralised
control centers and connected assets, this is achieved by connecting the RTU to a IEC 101 master, which then connects to a IEC 104 server, facilitating communication via TCP/IP \cite{6672100}. It is this IEC 104
modem which provides the network communications with the DSO SCADA system, as well as a number of interfaces to connect external devices, including serial and Ethernet connectivity.
It is these physical ports that will enable connection with the edge device, allow for the capture and parsing of the switch status messages without interfering with the existing stream to the DSOs SCADA system.
\\
\\
Within the context of the FLISR Edge algorithms operation then, the edge device will act as a relay to the centralised FLSIR algorithm when network connectivity is available, by passing switch status messages
to the centralised algorithm via MQTT (Message Queuing Telemetry Transport) a lightweight, publish/subscribe bi-directional communication protocol designed for resource constrained devices and/or remote locations \cite{7123563}.
To enable this interaction, an MQTT broker is deployed on the target system (e.g., DSO systems) which connects to the centralised FLISR algorithm, with the edge device exposing an outward MQTT interface.
Control messages generated by the centralised FLISR algorithm are captured by the edge device and relayed to the RTU, in the event that remote communication is unavailable,
the edge device will instead leverage the local FLISR algorithm and pass the generated control message to the RTU through the IEC 104 connection, outlined in \hyperref[fig:rtu_interaction]{Figure \ref{fig:rtu_interaction}}.

\begin{figure}[h]
  \centering
  \includegraphics[width=0.65\textwidth]{fig_rtu_interaction}
  \caption{\label{fig:rtu_interaction}Edge device interaction with switch hardware, DSO system \& centralised FLISR algorithm}
\end{figure}
\clearpage

\noindent Switch status and control messages utilised to define actions to take (i.e., opening/closing of the circuit breaker) contain a number of data points, an example is given below in JSON format.
\\
\begin{lstlisting}[language=json]
  {"LinkAdd": 1,"ASDU_CA": 51908,"TypeID": 30,"IOA": 4,"Value": "0(OFF)"}
\end{lstlisting}

\noindent The key properties in this data are, ``ASDU\_CA'' which is the identifier used by the DSO SCADA system to link the switch to its SCADA identifier, ``IOA'' (Information Object Address) which defines the specific input/output
referenced on the switch, for example, the status of the power supply, and finally ``Value'' which is given as a string value, ``0(OFF)'' mapping to ``False'', with a value of ``1(ON)'' mapping to ``True''.
While LBFM and Recloser switches both follow this convention, the IOAs which map to FLISR specific outputs are distinct as per manufacturer specifications.
\\
As stated in \hyperref[sec:flisr]{Section 2.2} a switch situated on a damaged line will pass an FPI and LVI status message, whereas a switch which has simply lost power, but not faulted, will pass an LVI status.
An FPI on it's own, would not necessarily indicate a fault, for example, a tree branch intermittently knocking against a wire on a gusty day, may result in a FPI status each occurrence,
but would not affect supply and be considered a ``Transient'' fault. As a general rule, any switch which gives an FPI and LVI signal, due to physical damage, should be opened, and therefore the damaged line or lines are no longer ``live''.
The FPI and LVI message format for both a LBFM and Recloser switch are given below in addition to the corresponding control message to open or close the circuit breaker.

\begin{itemize}
  \item \textbf{LBFM switch FPI, LVI \& control structure}
\end{itemize}

\noindent The FPI value for a LBFM switch map to IOA 4, meaning a ``Fault Current'' has been detected.
\\
\begin{lstlisting}[language=json]
  {"LinkAdd": 1,"ASDU_CA": 51908,"TypeID": 30,"IOA": 4,"Value": "1(ON)"}
\end{lstlisting}

\noindent The LVI value maps to IOA 2, which specifies ``AC Supply Fail''.
\\
\begin{lstlisting}[language=json]
  {"LinkAdd": 1,"ASDU_CA": 51908,"TypeID": 30,"IOA": 2,"Value": "1(ON)"}
\end{lstlisting}

\noindent An IOA value of 8 is used to open (``1(ON)'') or close (''0(OFF)'') the circuit breaker.
\\
\begin{lstlisting}[language=json]
  {"LinkAdd": 1,"ASDU_CA": 51908,"TypeID": 30,"IOA": 8,"Value": "1(ON)"}
\end{lstlisting}

\begin{itemize}
  \item \textbf{Recloser switch FPI, LVI \& control structure}
\end{itemize}

\noindent The FPI value for a Recloser switch maps to an IOA of 27, specifying a ``End of Protection Sequence''.
\\
\begin{lstlisting}[language=json]
  {"LinkAdd": 1,"ASDU_CA": 23089,"TypeID": 30,"IOA": 27,"Value": "1(ON)"}
\end{lstlisting}

\noindent The LVI value maps to IOA 34 ``Controller Locked Out''.
\\
\begin{lstlisting}[language=json]
  {"LinkAdd": 1,"ASDU_CA": 23089,"TypeID": 30,"IOA": 34,"Value": "1(ON)"}
\end{lstlisting}

\noindent Finally, IOA 4096 is utilised to open (``1(ON)'') or close (''0(OFF)'') the circuit breaker.
\\
\begin{lstlisting}[language=json]
  {"LinkAdd": 1,"ASDU_CA": 23089,"TypeID": 30,"IOA": 4096,"Value": "1(ON)"}
\end{lstlisting}

\noindent In operation then, both FLISR algorithms are designed to capture these switch status messages, and based on what is received (FPI and/or LVI) and the topology data related to that pole,
whether it is a LBFM/Recloser switch and/or a normally open point, determine the appropriate control message to return for a particular on pole circuit breaker. Similar to the simulation platform
services, both the centralised and distributed FLISR algorithms are implemented as lightweight containerised services, and are built on the ``golang:1.20-alpine'' Docker image. Both algorithms are defined in
the Go language, due to its suitability for building lightweight microservices, performance and efficiency. A study by Dymora et al., remarked that for processing small amounts of data,
Go surpassed comparable languages in terms of performance, CPU and memory utilisation \cite{app10238521}, making it an ideal candidate for implementation of the FLISR algorithms where reaction speed is paramount.
\\
As described in \hyperref[sec:objectives]{Section 1.3} both FLISR algorithms differ in their implementation, the centralised FLISR algorithm is deployed remotely and has access to the complete grid topology
hosted within an attached graph database. As such, this implementation of the algorithm will act on a fault event based on the status of all switches on an affected line. Conversely, the distributed FLISR algorithm
operates locally, and has access only to the data associated with the specific switch it is monitoring, \hyperref[fig:flisr_logic]{Figure \ref{fig:flisr_logic}} describes the basic logic of the FLISR algorithm operation
when a fault event has been detected.

\begin{figure}[h]
  \centering
  \includegraphics[width=0.5\textwidth]{fig_basic_flisr_logic}
  \caption{\label{fig:flisr_logic}Core FLISR Edge logic}
\end{figure}


% subsubsection flisr_edge (end)
